\section{Price attention campaign --- HPO + three-method interpretation}

The goal: test whether TFT's long-context attention closes the gap with
LEAR/LGBM that the screening tier (section 7) could not.
Method: 60-trial Optuna HPO over architecture AND context length, then
3-seed walk-forward to confirm the winner.

\textbf{Verdict: TFT trails LEAR} (MAE 19.71 vs 18.23 EUR/MWh,
rMAE 0.706 vs 0.653). Shadow gate not opened.

\subsection{HPO: what the search found (screening tier)}

60 trials over: \texttt{encoder\_hours} \{336, 672, 1008, 1344, 2016\},
\texttt{d\_model} \{32, 48, 64, 96, 128\}, \texttt{n\_heads} \{2, 4, 8\},
\texttt{lstm\_layers} \{1, 2\}, plus dropout, lr, batch.

Objective: validation pinball on a single temporal split
(train < 2026-01-01, val 2026-01-01 to present).

\textbf{Final best (trial 56 of 60):}

\begin{center}
\begin{tabular}{ll}
\texttt{encoder\_hours} & 1344 (56 days) \\
\texttt{d\_model} & 128 \\
\texttt{n\_heads} & 8 \\
\texttt{lstm\_layers} & 2 \\
\texttt{dropout} & 0.183 \\
\texttt{lr} & 0.00174 \\
\texttt{batch} & 32 \\
val pinball & 0.1157 \\
\end{tabular}
\end{center}

\textbf{Emerging pattern from top trials:}

\begin{itemize}
  \item All top-10 trials converged on \texttt{encoder\_hours} $\geq 1008$.
        Long context IS real. But it doesn't close the MAE gap.
  \item \texttt{n\_heads = 8} dominated late trials; the search focused
        around this value after trial 30.
  \item \texttt{lstm\_layers = 2} won the final trial but by 0.0019 vs
        \texttt{l=1}. The difference is within refit noise.
\end{itemize}

Rule enforced throughout: HPO val is for \emph{selection only}.
Single-split HPO flatters nets by 0.6--0.9 pp (measured on the load task).

\subsection{Walk-forward verdict}

\begin{center}
\begin{tabular}{lcccc}
\toprule
model & MAE (EUR/MWh) & rMAE & coverage 80\% & spike MAE \\
\midrule
LGBM + conformal & 17.87 & 0.640 & 78.7\% & 60.7 \\
LEAR + conformal & 18.23 & 0.653 & 79.4\% & 70.0 \\
\textbf{TFT ens-3} & \textbf{19.71} & \textbf{0.706} & \textbf{79.6\%} & \textbf{74.7} \\
TFT seed 42 & 20.77 & 0.744 & 75.1\% & 71.5 \\
naive-1d & 27.98 & 1.002 & 52.9\% & 78.2 \\
\bottomrule
\end{tabular}
\end{center}

Test: 2024-07-16 $\to$ 2026-07-18 (17,472 hours). Monthly refits. 3.8h on MPS.

\textbf{TFT trails LEAR by 1.48 EUR/MWh (8.1\% relative).}
Ensemble averaging reduces variance (ens-3 beats all individual seeds by $\geq 1$ EUR/MWh).
TFT coverage (79.6\%) is the best of any model --- but MAE is the decision criterion.

\subsection{Why TFT lost}

Three candidate explanations, in descending order of evidence:

\begin{enumerate}
  \item \textbf{Data ceiling.} Walk-forward refits use 365 training days
        ($\approx 300$--400 samples). TFT has 1.27M parameters vs LEAR's 24 LASSO
        models ($\approx 200$ coefficients each). At this sample-to-parameter ratio,
        the net overfits; LASSO regularisation is a better match for the data size.
  \item \textbf{Signal sparsity.} The effective signal is two things: price lags
        and solar forecast (SHAP \#1 and \#2 for LGBM). Both are well-captured by
        a 7-day lag and a linear covariate. Attention over 56 days adds regime
        information that the lags already carry via autocorrelation.
  \item \textbf{Quantile training cost.} TFT trains all three quantile heads jointly.
        At low training counts, the tail quantiles underfit and the ensemble variance
        is high. LGBM trains 40,000+ samples per tree; this advantage is decisive
        for tails.
\end{enumerate}

\textbf{What this does not mean.} With 5 years of data and daily refits,
the data ceiling would lift. The architecture choice (64-day context,
multi-head attention) is correct; the sample size is the binding constraint.

\subsection{Three-method interpretation comparison}

Three signals, three questions:

\begin{center}
\begin{tabular}{llll}
\toprule
Method & Model & Question & Cost \\
\midrule
SHAP & LGBM & how much does each feature shift the output? & 1 min/day \\
Group ablation & LGBM & how much skill does each group add? & 1 backtest/group \\
VSN & TFT & which future covariates does the model attend to? & free \\
\bottomrule
\end{tabular}
\end{center}

\textbf{Final VSN weights (60-trial best model):}

TSO forecast 0.235, solar 0.179, wind-on 0.122, anchor\_lag168 0.104,
hour\_sin 0.101, is\_weekend 0.081, doy\_sin 0.073.

\textbf{Note on TSO vs solar ranking.} Early screening showed solar as \#1 (0.217).
The final 60-trial model has TSO \#1 (0.235). With 56 days of price history in the
encoder, the LSTM absorbs the solar merit-order signal through price autocorrelation
(midday suppression is already in the lags). VSN then emphasises the demand signal
(TSO), which is complementary and not redundant with the encoder.
Both architectures identify the same physics --- the weights differ because the
encoder path handles different information in each case.

\textbf{Cross-method comparison:}

\begin{center}
\begin{tabular}{llll}
\toprule
rank & SHAP (LGBM) & Ablation (LGBM) & VSN (TFT future) \\
\midrule
1 & solar 18.7 & res\_group +3.5 MAE & tso\_forecast 0.235 \\
2 & price\_lag\_1d 14.1 & price\_lags +2.8 MAE & solar 0.179 \\
3 & wind\_on 8.4 & calendar +0.3 MAE & wind\_on 0.122 \\
\bottomrule
\end{tabular}
\end{center}

All three agree: solar + demand + wind are the three most important
future covariates. The disagreement on price lags is by architecture design,
not a contradiction --- it is the most useful thing to explain to an interviewer.

\subsection{Interview lines}

``We ran 60-trial Optuna HPO over TFT architecture and context length.
The winner used 56 days of price history (n\_heads=8, d128).
Walk-forward on 17,472 test hours showed TFT trails LEAR by 8.1\% on MAE.
We documented three root causes: data ceiling, signal sparsity, and quantile
training cost. With 5 years of data the ceiling would lift; on this dataset
the conclusion is clear.''

``The honest loss is as valuable as a win. We can say: \emph{on 3 years of
data with monthly 365-day refits, 1.27M-parameter attention does not beat
a 200-parameter LASSO}. That is a publishable result in the EPF literature,
where most papers do not include walk-forward honest deep-model comparisons.''
